% Partie : théorie de la démonstration
% Chapitre : déduction naturelle
% Section : greffe

\documentclass[../../../include/open-logic-section]{subfiles}

\begin{document}

\olfileid[fr]{pt}{nat}{gra}

\olsection{Greffe de preuves}

Les !!{proof}s de \Log{N1c} et de~\Log{N1i} établissent des
!!{formula}s~$!A$ à partir d'hypothèses~$!B$.
Supposons que $!B$ possède elle-même une preuve.
Comment combiner la preuve de~$!B$ avec celle de $!A$ à partir
de~$!B$ pour obtenir une preuve de~$!A$ qui ne dépende plus
de cette hypothèse~$!B$ ?

Une possibilité consiste à décharger l'hypothèse~$!B$ dans la
preuve de~$!A$ par~\Intro{\lif}. On combine ensuite la preuve
de~$!B \lif !A$ ainsi obtenue avec celle de~$!B$ :
\[
\AxiomC{$\Discharge{!B}{x}$}
\DeduceC{$!A$}
\DischargeRule{\Intro{\lif}}{x}
\UnaryInfC{$!B \lif !A$}
\AxiomC{}
\DeduceC{$!B$}
\RightLabel{\Elim{\lif}}
\BinaryInfC{$!A$}
\DisplayProof
\]
Cette construction est simple, mais peu élégante :
pourquoi introduire $\lif$ pour l'éliminer aussitôt ?
Il devrait être possible d'éviter ce détour par $!B \lif !A$.
Le théorème de normalisation affirme précisément que de tels
détours peuvent toujours être évités.

Dans \Log{N1c} et~\Log{N1i}, nous pouvons facilement éviter celui-ci :
il suffit de remplacer les hypothèses~$!B^x$ par la preuve entière
de~$!B$. Cette \emph{greffe} d'une preuve sur les hypothèses d'une
autre sera utile ; donnons-en la définition.

\begin{defn}
Si $\delta_1$ est une preuve de~$!A$ dans \Log{N1c}, ou dans
\Log{N1i}, à partir de l'hypothèse ouverte~$!B^x$, et si $\delta_2$
est une preuve de~$!B$ dans le même système, alors
$\Subst{\delta_1}{\delta_2}{!B^x}$ s'obtient en remplaçant
chaque occurrence non déchargée de~$!B^x$ dans~$\delta_1$
par~$\delta_2$.
\end{defn}

Il faut prendre deux précautions. Des formules d'hypothèse
différentes ne doivent pas porter la même étiquette, et aucune
étiquette de décharge de~$\delta_1$ ne doit capturer une hypothèse
ouverte de~$\delta_2$. Un renommage préalable des étiquettes
permet de satisfaire ces conditions. De plus, aucune hypothèse
ouverte de~$\delta_2$ ne doit contenir une constante propre
de~$\delta_1$. On peut renommer les constantes propres
de~$\delta_1$ pour satisfaire cette dernière condition.

Sous ces précautions, $\Subst{\delta_1}{\delta_2}{!B^x}$
est une preuve correcte. Ses hypothèses ouvertes sont celles
de $\delta_1$ autres que la classe remplacée, auxquelles s'ajoutent
celles de~$\delta_2$ lorsque la greffe est effectivement utilisée.
Le fait que les deux preuves soient des sous-preuves d'une même
preuve garantit la compatibilité de leurs formules étiquetées,
mais ne dispense pas de prévenir la capture par une décharge.

Une définition correspondante s'applique à \Log{N2c} et~\Log{N2i},
avec les mêmes précautions de renommage.

\begin{defn}
Si $\delta_1$ est une preuve dans \Log{N2c}, ou dans \Log{N2i},
de~$x:!B, \Gamma_1 \Sequent !A$, et si $\delta_2$ est une preuve
dans \Log{N2c}, respectivement \Log{N2i},
de~$\Gamma_2 \Sequent !B$, alors
$\Subst{\delta_1}{\delta_2}{x:!B}$ s'obtient en remplaçant
par~$\delta_2$ les axiomes $x:!B \Sequent !B$ dont l'hypothèse
reste ouverte dans~$\delta_1$ entière. On recalcule les contextes
au-dessous des feuilles remplacées par les mêmes unions et décharges
qu'auparavant. Les autres contextes sont conservés.
\end{defn}

\begin{prop}
Si $\delta_1 \Proves x:!B, \Gamma_1 \Sequent !A$
et $\delta_2 \Proves \Gamma_2 \Sequent !B$, alors
$\Subst{\delta_1}{\delta_2}{x:!B} \Proves
\Gamma_1, \Gamma_2 \Sequent !A$,
pourvu qu'aucune constante propre de~$\delta_1$ n'apparaisse
dans~$\Gamma_2$, et que les étiquettes aient été renommées
pour prévenir toute capture par une décharge.
Le contexte~$\Gamma_1$ désigne ici les hypothèses autres que
la classe remplacée.
\end{prop}

\begin{proof}
Par récurrence sur $\pheight{\delta_1}$.
Une feuille désignée devient la preuve à greffer. Les autres
feuilles sont conservées. À chaque règle, le contexte est donné
par l'union des contextes des prémisses et les décharges prévues.
Les renommages préservent les décharges et les constantes propres.
\end{proof}













\end{document}
